% Options for packages loaded elsewhere
\PassOptionsToPackage{unicode}{hyperref}
\PassOptionsToPackage{hyphens}{url}
\documentclass[
  11pt,
  a4paper,
]{article}
\usepackage{xcolor}
\usepackage[margin=24mm]{geometry}
\usepackage{amsmath,amssymb}
\setcounter{secnumdepth}{5}
\usepackage{iftex}
\ifPDFTeX
  \usepackage[T1]{fontenc}
  \usepackage[utf8]{inputenc}
  \usepackage{textcomp} % provide euro and other symbols
\else % if luatex or xetex
  \usepackage{unicode-math} % this also loads fontspec
  \defaultfontfeatures{Scale=MatchLowercase}
  \defaultfontfeatures[\rmfamily]{Ligatures=TeX,Scale=1}
\fi
\usepackage{lmodern}
\ifPDFTeX\else
  % xetex/luatex font selection
\fi
% Use upquote if available, for straight quotes in verbatim environments
\IfFileExists{upquote.sty}{\usepackage{upquote}}{}
\IfFileExists{microtype.sty}{% use microtype if available
  \usepackage[]{microtype}
  \UseMicrotypeSet[protrusion]{basicmath} % disable protrusion for tt fonts
}{}
\makeatletter
\@ifundefined{KOMAClassName}{% if non-KOMA class
  \IfFileExists{parskip.sty}{%
    \usepackage{parskip}
  }{% else
    \setlength{\parindent}{0pt}
    \setlength{\parskip}{6pt plus 2pt minus 1pt}}
}{% if KOMA class
  \KOMAoptions{parskip=half}}
\makeatother
\usepackage{longtable,booktabs,array}
\newcounter{none} % for unnumbered tables
\usepackage{calc} % for calculating minipage widths
% Correct order of tables after \paragraph or \subparagraph
\usepackage{etoolbox}
\makeatletter
\patchcmd\longtable{\par}{\if@noskipsec\mbox{}\fi\par}{}{}
\makeatother
% Allow footnotes in longtable head/foot
\IfFileExists{footnotehyper.sty}{\usepackage{footnotehyper}}{\usepackage{footnote}}
\makesavenoteenv{longtable}
\setlength{\emergencystretch}{3em} % prevent overfull lines
\providecommand{\tightlist}{%
  \setlength{\itemsep}{0pt}\setlength{\parskip}{0pt}}
\usepackage{bookmark}
\IfFileExists{xurl.sty}{\usepackage{xurl}}{} % add URL line breaks if available
\urlstyle{same}
\hypersetup{
  pdftitle={DBA Ring Lattice: Hamiltonian{,} KV Benchmark{,} and AI Design},
  hidelinks,
  pdfcreator={LaTeX via pandoc}}

\title{DBA Ring Lattice: Hamiltonian, KV Benchmark, and AI Design}
\author{}
\date{1 November 2026}

\begin{document}
\maketitle

All file paths below are relative to the project root, one level above
this TeX directory.

\section{Synthetic ring Hamiltonian}\label{synthetic-ring-hamiltonian}

This is a complete quadratic Hamiltonian for a specified synthetic
linear machine. It is not an exact nonlinear magnetic-field Hamiltonian.
Hard edges, paraxial motion, first-order momentum deviations and a
static reference energy define the model. Fringe fields, chromatic
nonlinearities, RF, radiation and errors are excluded. The DBA-like
geometry is deliberately non-achromatic.

\subsection{Canonical convention}\label{canonical-convention}

Independent variable: reference arc length s {[}m{]}. Canonical state:

\[
\begin{aligned}
\mathbf u&=(x,p_x,y,p_y,\zeta,\eta),\\
p_x&=P_x/p_0,\qquad p_y=P_y/p_0,\\
\zeta&=\beta_0c\left(\frac{s}{\beta_0c}-t\right)\quad[\mathrm m],\\
\eta&=\frac{E-E_0}{\beta_0^2E_0}\quad\text{(dimensionless)}.
\end{aligned}
\]

P denotes canonical transverse momentum; choose the usual ideal magnet
gauge with zero transverse vector potential. Slopes equal normalized
transverse momenta to linear order. eta equals relative momentum
deviation dp/p0 only to first order. beta0 and gamma0 refer to 2.5 MeV
protons. This is not OPALX's momentum normalization by mc; the
corresponding on-momentum conversion is P\_x/(mc)=beta0\emph{gamma0}px.

Positive \(\zeta\) means early arrival. Local x points outward in a
bend, y is vertical. Positive curvature h bends the reference towards
negative lab X at the entrance. The lab tangent in the X,Z plane is
(sin(theta),cos(theta)), with theta'=-h. Entry and exit placements in
\texttt{model/machine.json} use this convention, matching the copied
baseline deck.

The scaled, dimensionless Hamiltonian is

\[
\begin{aligned}
H_{\mathrm{ext}}(\mathbf u;s)&=\frac{p_x^2+p_y^2}{2}
 +\frac{k_x(s)x^2+k_y(s)y^2}{2}-h(s)x\eta+\frac{\eta^2}{2\gamma_0^2},\\
x'&=p_x,&p_x'&=-k_xx+h\eta,\\
y'&=p_y,&p_y'&=-k_yy,\\
\zeta'&=-hx+\eta/\gamma_0^2,&\eta'&=0.
\end{aligned}
\]

h has units m\^{}-1 and kx,ky have units m\^{}-2. H is piecewise
constant in s through its coefficients; coordinates and canonical
momenta remain continuous at hard edges. No thin edge kicks are applied.
Use \texttt{u\textquotesingle{}=J*G*u}, with G the Hessian of H and
J=diag(J2,J2,J2). The exact element map is exp(L\emph{J}G); ordered cell
and ring maps are M\_cell=M\_B2 M\_D2 M\_Q M\_D1 M\_B1 and
M\_ring=M\_cell\^{}6. An element's autonomous H is conserved within that
element; the piecewise s-dependent H is not a global conserved scalar.
The exact maps are symplectic, including dispersion and longitudinal
coupling.

\subsection{Cell definition}\label{cell-definition}

{\def\LTcaptype{none} % do not increment counter
\begin{longtable}[]{@{}lrrrr@{}}
\toprule\noalign{}
Element & L {[}m{]} & h {[}m\^{}-1{]} & kx {[}m\^{}-2{]} & ky
{[}m\^{}-2{]} \\
\midrule\noalign{}
\endhead
\bottomrule\noalign{}
\endlastfoot
B1 & pi/3 & 0.5 & 0.25 & 0 \\
D1 & 1 & 0 & 0 & 0 \\
Q & 0.2 & 0 & -2.84 & +2.84 \\
D2 & 1 & 0 & 0 & 0 \\
B2 & pi/3 & 0.5 & 0.25 & 0 \\
\end{longtable}
}

Six cells give 30 elements, twelve 30-degree bends, total bending 2*pi,
and circumference 25.76637061436 m. The proton OPALX deck uses K1=-2.84
to implement these physical quadrupole coefficients. The electron
historical benchmark uses the opposite OPALX K1 sign. Do not infer
physical focusing from an OPALX K1 without its particle/charge
convention.

Bare full tunes are Qx=0.770107569133 and Qy=1.689536397020. The PR's
principal one-turn eigenphases 0.229892431 and 0.310463603 are
compatible with these values; they are not the full positive tunes. The
model generator chooses cell phase in (0,2*pi) using the sign of M12 and
multiplies by six. This branch is verified by the independent KV phase
integral at the baseline. Local optimization stays near that branch; a
general design with more than one oscillation per cell must unwrap phase
along s independently. The periodic entrance dispersion is approximately
9.337460094 m and zero slope. The machine has longitudinal drift, not
longitudinal confinement.

\subsection{Matched KV beam}\label{matched-kv-beam}

For the solver comparison use eta=0 and zero momentum spread. An
initially uniform coasting beam remains longitudinally uniform in the
ideal model. Let a,b be edge radii {[}m{]}, a=2\emph{sigma\_x,
b=2}sigma\_y. Geometric edge emittances \(\mathcal E_x,\mathcal E_y\)
(Ex,Ey in the code) are four times geometric rms emittances; the
baseline has rms emittances 3e-6 m rad. Current I is positive for the
proton beam.

\[
\begin{aligned}
I_A&=\frac{4\pi\epsilon_0m_pc^3}{e},&K&=\frac{2I}{I_A\beta_0^3\gamma_0^3},\\
a''+k_xa-\frac{\mathcal E_x^2}{a^3}-\frac{2K}{a+b}&=0,\\
b''+k_yb-\frac{\mathcal E_y^2}{b^3}-\frac{2K}{a+b}&=0.
\end{aligned}
\]

For particles inside the uniform KV ellipse the collective transverse
potential adds

\[
\begin{aligned}
H_{\mathrm{sc}}&=-\frac{K}{a+b}\left(\frac{x^2}{a}+\frac{y^2}{b}\right),\\
k_{x,\mathrm{eff}}&=k_x-\frac{2K}{a(a+b)},\\
k_{y,\mathrm{eff}}&=k_y-\frac{2K}{b(b+a)}.
\end{aligned}
\]

This is an interior test-particle mean-field Hamiltonian with
self-consistently matched envelopes. It is not a quadratic potential
outside the ellipse, not a sum of finite-particle pair interactions, and
not a six-dimensional collective model with longitudinal space charge.
Self-consistency is supplied by the periodic envelope equations, not by
holding the zero-current beam sizes fixed as current rises. Additive
functions of s do not change transverse motion.

The KV distribution is a uniform shell in normalized four-dimensional
transverse phase space. A filled 4D ball is a different distribution.
Its invariant is

\[
\left(\frac{x}{a}\right)^2+
\left(\frac{ap_x-a'x}{\mathcal E_x}\right)^2+
\left(\frac{y}{b}\right)^2+
\left(\frac{bp_y-b'y}{\mathcal E_y}\right)^2=1.
\]

Full positive incoherent tunes follow Qx=6/(2\emph{pi)}integral\_cell
Ex/a\^{}2 ds, and similarly for y. The self-field is analytic; periodic
envelope matching and phase integration are numerical. Distinguish these
particle tunes from centroid tunes and envelope-mode frequencies. Finite
amplitudes/curvature in an OPALX field tracker can depart from this
paraxial locally straight space-charge reference; mesh/time convergence
alone does not remove that model difference.

\subsection{References}\label{references}

\begin{itemize}
\tightlist
\item
  \href{https://indico.cern.ch/event/221441/contributions/1524769/attachments/360572/502106/SpaceCharge_Modelling.pdf}{CERN
  CAS: modelling with space charge, C. Prior}, for KV envelope and
  perveance conventions.
\item
  \href{https://s3.cern.ch/inspire-prod-files-7/74c88492a8d6a790563d145198c4813d}{Particle
  tracking in circular accelerators using the exact Hamiltonian}, for
  canonical accelerator Hamiltonian context. The explicit quadratic
  convention and signs used here are defined above and checked against
  closed-form sector maps.
\end{itemize}

\section{KV reference for FFT2D5}\label{kv-reference-for-fft2d5}

The first deliverable is an independent reference, not yet an end-to-end
OPALX FFT2D5 validation. \texttt{tune-shift-ring/analysis/run\_model.py}
regenerates immutable current scans, charge normalization, periodic
envelopes, tune predictions, stability diagnostics and local
electric-field gradients. \texttt{fft2d5\_reference.json} in each new
result directory records numerical comparison targets.

\subsection{Physical comparison}\label{physical-comparison}

Use 2.5 MeV protons, 3 micrometre geometric rms emittances in both
planes, zero momentum spread, no RF, no pipe images and open transverse
boundaries. Match the beam separately at each charge. The coasting
reference is uniform longitudinally. In the current OPALX
implementation, FFT2D5 uses lab-frame line charge and open transverse
slice solves. Its longitudinal-field mode must be None for a controlled
purely transverse comparison (C++ enum
\texttt{FFT2D5LongitudinalFieldMode::None}; verify the input spelling
when building a run deck).

Charge normalization is determined by line density, not by an arbitrary
bunch charge:

\[
\begin{aligned}
\lambda&=\frac{I}{\beta_0c},&N_{\mathrm{ring}}&=\frac{\lambda C}{e},\\
Q_{\mathrm{ring}}&=\lambda C,&Q_{\mathrm{segment}}&=\lambda L_{\mathrm{segment}},\\
q_{\mathrm{macro}}&=\frac{Q_{\mathrm{segment}}}{N_{\mathrm{macro}}}.
\end{aligned}
\]

A full coasting ring needs a spatial distribution all around the closed
path with the appropriate local matched covariance at each s. A finite
segment requires verified uniform interior slices and appropriate
longitudinal handling. A zero-length observation section
(\texttt{kv-section.txt}) cannot represent a charged coasting beam. Do
not replace the coasting model with an all-open isolated ellipsoid.
Full-ring longitudinal wrapping and the path seam must be verified
before claiming a ring comparison.

At each cell position the analytic lab electric field inside the ellipse
is

\[
\begin{aligned}
E_x&=\frac{\lambda x}{\pi\epsilon_0a(a+b)},&
E_y&=\frac{\lambda y}{\pi\epsilon_0b(a+b)},\\
\mathbf B&=\frac{\beta_0}{c}\,\hat{\mathbf t}\times\mathbf E,&
\mathbf F_\perp&=\frac{e\mathbf E_\perp}{\gamma_0^2}.
\end{aligned}
\]

The electric field already uses lab normalization; do not add another
gamma boost. The force reduction is from the magnetic field.
Equivalently,
e\emph{(E\_x/x)/(p0}beta0\emph{c}gamma0\^{}2)=2*K/{[}a*(a+b){]}, which
ties a field check directly to the Hamiltonian.

\subsection{Baseline operating points}\label{baseline-operating-points}

{\def\LTcaptype{none} % do not increment counter
\begin{longtable}[]{@{}
  >{\raggedleft\arraybackslash}p{(\linewidth - 10\tabcolsep) * \real{0.1739}}
  >{\raggedleft\arraybackslash}p{(\linewidth - 10\tabcolsep) * \real{0.1739}}
  >{\raggedleft\arraybackslash}p{(\linewidth - 10\tabcolsep) * \real{0.1739}}
  >{\raggedleft\arraybackslash}p{(\linewidth - 10\tabcolsep) * \real{0.1739}}
  >{\raggedleft\arraybackslash}p{(\linewidth - 10\tabcolsep) * \real{0.1739}}
  >{\raggedright\arraybackslash}p{(\linewidth - 10\tabcolsep) * \real{0.1304}}@{}}
\toprule\noalign{}
\begin{minipage}[b]{\linewidth}\raggedleft
I {[}mA{]}
\end{minipage} & \begin{minipage}[b]{\linewidth}\raggedleft
Qx
\end{minipage} & \begin{minipage}[b]{\linewidth}\raggedleft
Qy
\end{minipage} & \begin{minipage}[b]{\linewidth}\raggedleft
Delta Qx
\end{minipage} & \begin{minipage}[b]{\linewidth}\raggedleft
Delta Qy
\end{minipage} & \begin{minipage}[b]{\linewidth}\raggedright
Role
\end{minipage} \\
\midrule\noalign{}
\endhead
\bottomrule\noalign{}
\endlastfoot
0 & 0.770108 & 1.689536 & 0 & 0 & Bare optics \\
4.748770 & 0.623940 & 1.589536 & -0.146167 & -0.1 & First collective
comparison \\
10.399311 & 0.484005 & 1.489536 & -0.286102 & -0.2 & Main
substantial-shift reference \\
17.402677 & 0.359252 & 1.389536 & -0.410855 & -0.3 & Envelope-unstable
stress case \\
\end{longtable}
}

At the main point Q\_ring=12.26826245 nC and N\_ring=7.65724714e10
protons. Entrance rms sizes are 6.042810 mm and 2.480302 mm. The last
operating point has an envelope multiplier about 1.069761 per cell and
is not a stable matched-beam benchmark.

\subsection{Validation sequence}\label{validation-sequence}

\begin{enumerate}
\def\labelenumi{\arabic{enumi}.}
\tightlist
\item
  Reconstruct the bare ring and establish convergence of the
  particle/reference orbit and zero-charge tunes in OPALX. Historical
  tracking reports describe old source revisions and do not establish
  current behavior.
\item
  At a fixed matched section compare FFT2D5 interior field gradients and
  net transverse force against the formulas above. Check charge
  deposition and line density as well as fields. Repeat at a curved
  section to expose path/coordinate errors.
\item
  Transport the correctly normalized matched source through one cell.
  Compare envelopes and projected moments against the reference at every
  element boundary.
\item
  Track stable coasting operating points over many turns. Measure
  individual-particle tunes with a consistent observation plane and
  unfolded branches, not the centroid or only a principal FFT peak.
  Include the zero-charge run at the same numerical settings.
\item
  Quantify sensitivity to timestep, transverse mesh spacing/domain,
  slice spacing, source particle count and spectral observation length.
  Keep physical parameters fixed and report both absolute tune error and
  error in Delta Q. Establish tolerances from this convergence evidence
  rather than choosing a pass threshold now.
\end{enumerate}

Parameter studies use one MPI rank. Dedicated single-/multi-rank
consistency checks belong to the eventual solver validation, not these
serial reference calculations. No OPALX or MPI solver validation has
been run for this standalone package yet.

\section{AI task: reconstruct, then
optimize}\label{ai-task-reconstruct-then-optimize}

Provide an AI with \texttt{model/cell.json} and the Hamiltonian section
of this document. Keep \texttt{model/machine.json} as the evaluator's
reference output if testing independent reconstruction.

\subsection{Stage 1: reconstruct the
machine}\label{stage-1-reconstruct-the-machine}

Read the declared canonical variables and quadratic Hamiltonian. Derive
the six Hamilton equations, generate the exact map of each element,
compose a cell in beam order, and repeat it six times. Integrate the
reference geometry and output all 30 placements, the cell/ring maps,
circumference, bend sum, dispersion, full tunes and symplectic
residuals. Do not infer full tunes from principal one-turn eigenphases
alone.

Use the proton sign convention. At zero current reproduce the
independent envelope phase integrals. Explain the distinction between a
full quadratic machine model and exact nonlinear field tracking.

The supplied baseline evaluator is \texttt{model/machine.py}; physics
tests are in \texttt{tests/test\_machine.py}. Baseline targets:
25.76637061436 m circumference, 2*pi total bend, full tunes
(0.770107569133,1.689536397020), and entrance dispersion (9.337460094
m,0). Tests use near-roundoff tolerances for exact maps and geometry and
2e-10 for comparison with numerically integrated envelope phases. These
tolerances are model checks, not PIC acceptance criteria.

\subsection{Stage 2: optimize the cell}\label{stage-2-optimize-the-cell}

For a first explicit demonstration, choose full tune targets (0.8,1.7).
These are illustrative targets, not a requested operating machine or an
optimal physics choice. Allow the physical quadrupole coefficient and
both equal drift lengths to vary. Hold particle, energy, bend
radii/angles, quadrupole length and cell count fixed. Require positive
lengths, transverse cell stability, symplectic maps, and closed
reference geometry. Circumference is allowed to change. If circumference
is fixed, these two knobs no longer generally permit two arbitrary tune
targets; add a justified degree of freedom or change the objective.

Export the new cell specification and full machine, plus a KV
configuration derived from exactly the same element coefficients.
Rematch the KV beam and recompute the tune-shift-versus-charge curve for
each proposed design. Reject envelope-unstable points as stable working
references. Optimization of bare tunes alone does not optimize
collective stability or machine robustness.

\texttt{model/optimize.py} is a deterministic local two-knob reference
optimizer, useful for evaluating an AI's proposal. It is not evidence
that AI independently designed a machine. It deliberately fails when
local matching encounters an unstable/nonphysical cell or cannot
converge; it is not a global search over tune branches.

The worked result in \texttt{designs/demo-qx0.8-qy1.7/} reconstructs a
complete machine with full tunes (0.8,1.7). Its compatible
\texttt{kv-config.json} can be supplied to the analytic scan runner.
Compare proposals on achieved optics, valid geometry, collective
stability, reproducibility and stated tradeoffs, rather than
plausible-looking element lists.

\end{document}
